% esl-mpsoc-memory-hierarchy — an MPSoC tile array sharing one memory block over a pinned inter-tile bus
% THESIS: only the shared memory block crosses tile boundaries -- every PE stays local to its tile -- so the single inter-bus fan-in reveals the design's memory-boundedness at a glance
% sources: tsarilp-pub doc/common/tikzsetup.tex; cp-4856 esl-architecture-block PoC graft; cp-4892 esl-mpsoc-tile-array foreach generator layer
% constraints: ADR-0005 D3 §1b row: 1:1 node_naming correspondence with the contract component vocabulary; ADR-0005 D5: array scaling stays foreach-driven, no hard-coded node counts
\documentclass[border=8pt]{standalone}

% --- packages (mirror these in template.meta.json "requires") ---
\usepackage[dvipsnames]{xcolor}
\usepackage{tikz}
\usetikzlibrary{calc, backgrounds, fit, positioning, arrows.meta}

% --- palette: ESL domain-semantic tokens (contract §2, ADR-0005 D7) ---
\colorlet{core-accent}{Aquamarine}     % PE fill accent
\colorlet{intra-bus}{Orange}           % intra-tile interconnect
\colorlet{inter-bus}{NavyBlue}         % inter-tile interconnect
\colorlet{digitaldomain}{gray}         % shared memory block (digital storage)

\begin{document}
% ==== parameters (edit these; this is the \foreach generator's parameter surface —
% nothing below this block should ever be edited to change array shape/size) =====
\def\ntiles{3}          % number of tiles in the row
\def\pespertile{2}      % PEs per tile
\def\tilegap{0.7}       % horizontal gap between tile boundary boxes (cm)
\def\pegap{0.12}        % gap between PEs within a tile (cm)
\def\pesize{0.5}        % PE node minimum size (cm)
\def\memheight{1.6}     % vertical offset of the shared mem-block above the tile row (cm)
\def\tilelabelprefix{T} % tile label prefix
\def\memlabel{L2 \\ shared mem}  % shared mem-block label (use \\ for two lines)
% =========================================================================

\begin{tikzpicture}[
    >=Stealth,
    pe/.style={draw, thick, rectangle, minimum size=\pesize cm, inner sep=0pt,
      draw=core-accent!80!black, fill=core-accent!25, font=\tiny, align=center},
    tilebox/.style={draw=core-accent!70!black, rounded corners=3pt, inner sep=6pt},
    tilelabel/.style={font=\scriptsize\bfseries, text=core-accent!55!black},
    memblock/.style={draw=digitaldomain!80!black, thick, rectangle, rounded corners=2pt,
      fill=digitaldomain!18, minimum width=1.6cm, minimum height=0.7cm,
      font=\scriptsize, align=center},
    intrabus/.style={draw=intra-bus, thick},
    interbus/.style={draw=inter-bus, thick, dashed},
  ]

  % --- pitch: uniform spacing between successive tile origins ---
  \pgfmathsetmacro{\pestep}{\pesize+\pegap}
  \pgfmathsetmacro{\tilepitch}{\pespertile*\pestep+\tilegap}

  % ==== generator: tile row x { PE row, containment fit box } (contract:
  % tile ⊃ pe, primitive containment; §3a) ====
  \foreach \t in {1,...,\ntiles} {
    \pgfmathsetmacro{\ox}{(\t-1)*\tilepitch}
    \def\tilefitlist{}
    \foreach \k in {1,...,\pespertile} {
      \pgfmathsetmacro{\px}{\ox+(\k-1)*\pestep}
      \node[pe] (pe-\t-\k) at (\px,0) {\k};
      \xdef\tilefitlist{\tilefitlist(pe-\t-\k)}
    }
    \begin{scope}[on background layer]
      \node[tilebox, fit=\tilefitlist] (tile-\t) {};
    \end{scope}
    \node[tilelabel, anchor=north west]
      at ([xshift=2pt,yshift=-2pt]tile-\t.north west) {\tilelabelprefix\t};
    % intra-tile bus: solid chain across this tile's PEs (relative-placement
    % within the container)
    \ifnum\pespertile>1
      \foreach \k in {2,...,\pespertile} {
        \pgfmathtruncatemacro{\kprev}{\k-1}
        \draw[intrabus] (pe-\t-\kprev) -- (pe-\t-\k);
      }
    \fi
  }

  % ==== shared mem-block, pinned along one fixed row above every tile
  % (pinned-axis primitive, §3a: "interface/port row pinned on a tile edge") ====
  \coordinate (rowmid) at ($(tile-1.north)!0.5!(tile-\ntiles.north)$);
  \node[memblock, above=\memheight cm of rowmid] (mem) {\memlabel};
  % orthogonal trunk-and-branch routing (contract §3a #5 typed-routing:
  % "orthogonal, no stray crossings"; architecture realization: "orthogonal
  % intra/inter-bus at ports") -- a straight tile.north -- mem diagonal per
  % tile was this template's first draft, copying esl-architecture-block's
  % pre-fix defect; replaced with vertical tile stubs into one shared
  % horizontal rail, then one vertical drop into mem. \railmid sits at the
  % vertical midpoint between the tile-top row and mem, sharing mem's x (mem
  % is centered above rowmid, and railmid is derived from rowmid), so every
  % stub, the rail, and the drop meet with no diagonal segment.
  \coordinate (railmid) at ($(rowmid)!0.5!(mem)$);
  \foreach \t in {1,...,\ntiles} {
    \draw[interbus] (tile-\t.north) -- (tile-\t.north -| railmid);
  }
  \draw[interbus] (tile-1.north -| railmid) -- (tile-\ntiles.north -| railmid);
  \draw[interbus] (railmid) -- (mem);

\end{tikzpicture}
\end{document}
